\documentclass[11pt]{article}
\usepackage[margin=1in]{geometry}
\usepackage{amsmath,amssymb}
\usepackage{hyperref}
\usepackage{enumitem}
\title{Project Plan: From Radiation-Induced Defects to Device Degradation in Silicon}
\author{}
\date{}

\begin{document}
\maketitle

\section{Project Objective}
The goal of this project is to connect radiation transport, defect generation, semiconductor material degradation, and device-level performance loss in silicon. The project is organized so that a simple room-temperature model can be studied first, then extended to incorporate Geant4 outputs, and finally extended further to include cryogenic defect kinetics and persistent trapping behavior.

\section{Core Physics Story}
The main physics chain is:
\begin{center}
incident radiation $\rightarrow$ displacement-producing interactions $\rightarrow$ defects in silicon $\rightarrow$ changes in carrier lifetime, diffusion length, and leakage $\rightarrow$ measurable device degradation.
\end{center}

A cryogenic extension is then appended:
\begin{center}
defects created by radiation $\rightarrow$ temperature-dependent migration, trap emission, and annealing $\rightarrow$ frozen-in defect populations and slower recovery at low temperature.
\end{center}

\section{Project Modules}
\subsection{Module 1: Basic room-temperature MATLAB model}
This module provides a self-contained phenomenological model of radiation damage in silicon.

Representative relations include:
\begin{align}
N_D &= k_D \Phi, \\
\frac{1}{\tau_{\mathrm{eff}}} &= \frac{1}{\tau_0} + K_\tau N_D, \\
L &= \sqrt{D\tau_{\mathrm{eff}}}, \\
\Delta I &= \alpha \Phi V.
\end{align}

This module is used to generate first-pass plots of:
\begin{itemize}[leftmargin=2em]
    \item defect density vs fluence,
    \item carrier lifetime vs fluence,
    \item diffusion length vs fluence,
    \item leakage current vs fluence,
    \item charge collection trend vs fluence.
\end{itemize}

\subsection{Module 2: Geant4-coupled MATLAB model}
This module imports Geant4 outputs such as event summaries, depth-dependent energy deposition, and recoil-like proxies, then uses those as inputs to the degradation model.

Its purpose is to make the project more physically grounded by connecting radiation transport results to defect and device trends.

\subsection{Module 3: Cryogenic defect kinetics extension}
This extension models temperature-dependent defect behavior using simple Arrhenius forms:
\begin{align}
D_{\mathrm{def}}(T) &= D_0 \exp\left(-\frac{E_m}{k_B T}\right), \\
e(T) &= e_0 \exp\left(-\frac{E_a}{k_B T}\right), \\
R_{\mathrm{ann}}(T) &= R_0 \exp\left(-\frac{E_{\mathrm{ann}}}{k_B T}\right).
\end{align}

This module is used to show how migration, trap emission, and annealing all slow strongly as temperature decreases.

\subsection{Module 4: Persistent trapping and slow recovery extension}
This extension connects cryogenic defect kinetics to a concrete device-level consequence using a simple trapped-charge recovery model:
\begin{equation}
Q_{\mathrm{trap}}(t,T) = Q_0 \exp\left(-\frac{t}{\tau_{\mathrm{rec}}(T)}\right).
\end{equation}

This module illustrates that trapped charge and damage signatures can persist much longer at cryogenic temperature because recovery becomes very slow.

\section{Important Clarification: Development Order vs Run Order}
There are two different orderings in this project, and they should not be confused.

\subsection{Development order}
The development order refers to how the project was built conceptually and how the modules depend on one another in the broader design roadmap.

The recommended development order is:
\begin{enumerate}[leftmargin=2em]
    \item Build the core room-temperature MATLAB model.
    \item Add Geant4 coupling so transport outputs can feed the degradation model.
    \item Add the cryogenic defect kinetics extension.
    \item Add the persistent trapping and slow-recovery extension.
\end{enumerate}

This is the order that makes the project architecture clean and logically expandable.

\subsection{Practical run order}
The practical run order refers to the best order for actually executing the MATLAB scripts while learning the code and debugging the workflow.

The recommended practical run order is:
\begin{enumerate}[leftmargin=2em]
    \item Run \texttt{main\_full\_pipeline.m} first.
    \item Then run \texttt{main\_cryogenic\_extension.m}.
    \item Then run \texttt{main\_persistent\_trapping\_extension.m}.
    \item Run \texttt{main\_import\_geant4\_pipeline.m} only after Geant4 CSV outputs are available.
\end{enumerate}

This run order is easier because it does not require Geant4 to be working before the user can understand the room-temperature and cryogenic MATLAB physics.

\subsection{Why these two orders differ}
The development order is an integration roadmap. It explains how the full project was assembled conceptually.

The practical run order is a user workflow. It explains the easiest way to execute the scripts without being blocked by Geant4 setup.

In short:
\begin{itemize}[leftmargin=2em]
    \item development order = project architecture roadmap,
    \item run order = easiest execution and debugging roadmap.
\end{itemize}

\section{Recommended Use of Each Main MATLAB Script}
\subsection*{\texttt{main\_full\_pipeline.m}}
Use this first. It is the baseline room-temperature model and does not require Geant4 outputs.

\subsection*{\texttt{main\_import\_geant4\_pipeline.m}}
Use this only after Geant4 has produced required CSV files such as event summaries and depth-dependent energy deposition data.

\subsection*{\texttt{main\_cryogenic\_extension.m}}
Use this to study how defect migration, trap emission, and annealing change with temperature.

\subsection*{\texttt{main\_persistent\_trapping\_extension.m}}
Use this to study the device-level consequence of slow recovery and persistent trapping at cryogenic temperature.

\subsection*{Single-topic scripts}
Scripts such as \texttt{main\_defect\_model.m}, \texttt{main\_lifetime\_model.m}, \texttt{main\_diffusion\_length\_model.m}, and \texttt{main\_leakage\_model.m} are best treated as small standalone demos for inspecting one part of the model in isolation. They do not need to be run in sequence for the main project workflow.

\section{Minimal Working Path}
A good minimal working path for the project is:
\begin{enumerate}[leftmargin=2em]
    \item Run the room-temperature model.
    \item Run the cryogenic kinetics model.
    \item Run the persistent trapping extension.
    \item After Geant4 is working, connect the Geant4 outputs through the import pipeline.
\end{enumerate}

This path gives a clear progression from simple model to transport-coupled model while keeping the learning process manageable.

\section{Deliverables}
The project should ultimately produce:
\begin{itemize}[leftmargin=2em]
    \item a room-temperature defect-to-device degradation model,
    \item a Geant4-coupled transport-to-degradation workflow,
    \item a cryogenic defect kinetics extension,
    \item a persistent trapping and slow-recovery extension,
    \item plots suitable for technical presentations,
    \item supporting notes for both room-temperature and cryogenic interpretations.
\end{itemize}

\end{document}
